\hypertarget{ux5deux5eaux5d5ux5d3ux5d5ux5dcux5d5ux5d2ux5d9ux5d5ux5ea-ux5d4ux5e2ux5e8ux5dbux5d4}{%
\section{מתודולוגיות
הערכה}\label{ux5deux5eaux5d5ux5d3ux5d5ux5dcux5d5ux5d2ux5d9ux5d5ux5ea-ux5d4ux5e2ux5e8ux5dbux5d4}}

\hypertarget{ux5deux5d3ux5d3ux5d9-ux5d4ux5e2ux5e8ux5dbux5d4-ux5d1ux5e2ux5d9ux5d1ux5d5ux5d3-ux5e9ux5e4ux5d4-ux5d8ux5d1ux5e2ux5d9ux5ea}{%
\subsection{מדדי הערכה בעיבוד שפה
טבעית}\label{ux5deux5d3ux5d3ux5d9-ux5d4ux5e2ux5e8ux5dbux5d4-ux5d1ux5e2ux5d9ux5d1ux5d5ux5d3-ux5e9ux5e4ux5d4-ux5d8ux5d1ux5e2ux5d9ux5ea}}

הערכת מודלי שפה וארכיטקטורות טרנספורמר מהווה מרכיב מרכזי בתהליך הפיתוח
והמחקר. מדדי הערכה סטנדרטיים מאפשרים השוואה בין מודלים שונים ומעניקים
תמונה כמותית של ביצועיהם על משימות מוגדרות. בקרב המדדים הנפוצים ביותר
נמנים ציון BLEU למכונת תרגום, מדד F1 לסיווג ולזיהוי ישויות, ומדד הנבוכות
לבחינת מודלי שפה.

ציון BLEU (Bilingual Evaluation Understudy) נועד למדוד את הדמיון בין
תרגום מכונה לבין תרגום אנושי ייחוס. החישוב מבוסס על דיוק n-gram משוקלל
בתוספת עונש קיצור. נוסחת BLEU הבסיסית מוגדרת כדלקמן:

BLEU = BP * exp( sum\_\{n=1\}\^{}\{N\} w\_n * log(p\_n) )

כאשר BP הוא עונש הקיצור (Brevity Penalty), p\_n הוא דיוק ה-n-gram
המשוקלל, ו-w\_n הם משקלי הלוגריתמים. ציון BLEU נע בין 0 ל-1, כאשר ערך
גבוה יותר מצביע על קרבה רבה יותר לתרגום הייחוס.

מדד F1 משמש להערכת משימות סיווג כגון ניתוח סנטימנט וזיהוי ישויות שמות
(NER). הוא מחושב כממוצע הרמוני של דיוק (Precision) וזכירה (Recall):

F1 = 2 * (Precision * Recall) / (Precision + Recall)

מדד הנבוכות (Perplexity) משמש להערכת מודלי שפה הסתברותיים. ערך נמוך של
נבוכות מצביע על כך שהמודל מייחס הסתברות גבוהה יותר לרצפים שנדגמו מהשפה
האמיתית. הנבוכות מוגדרת כ:

PP(W) = P(w\_1, w\_2, \ldots, w\_N)\^{}\{-1/N\}

כאשר W הוא רצף המילים ו-N הוא אורכו. בפועל, מודלי שפה מודרניים מבוססי
טרנספורמר מושווים על בסיס נבוכות על קורפוסים סטנדרטיים.

\hypertarget{ux5deux5d3ux5d3ux5d9-ux5d4ux5e2ux5e8ux5dbux5d4-ux5dcux5deux5e9ux5d9ux5deux5d5ux5ea-ux5d4ux5d1ux5e0ux5ea-ux5e9ux5e4ux5d4}{%
\subsection{מדדי הערכה למשימות הבנת
שפה}\label{ux5deux5d3ux5d3ux5d9-ux5d4ux5e2ux5e8ux5dbux5d4-ux5dcux5deux5e9ux5d9ux5deux5d5ux5ea-ux5d4ux5d1ux5e0ux5ea-ux5e9ux5e4ux5d4}}

לצד המדדים הכמותיים הפשוטים, פותחו ערכות השוואה (benchmark suites)
מקיפות המאפשרות הערכה רב-ממדית של מודלי שפה. ערכת GLUE (General Language
Understanding Evaluation) כוללת תשע משימות שונות הבוחנות הבנת שפה, כגון
הסקה טקסטואלית, ניתוח סנטימנט ודמיון סמנטי בין משפטים. ערכת SuperGLUE
הורחבה מ-GLUE וכוללת משימות מאתגרות יותר, כגון מענה לשאלות מרובות הקשר
ומשימות הסקה לשונית.

ערכת SQuAD (Stanford Question Answering Dataset) מהווה סטנדרט מרכזי
למשימות מענה לשאלות. בגרסתה הראשונה, SQuAD 1.1, כל תשובה מופיעה כמקטע
ישיר מתוך הפסקה הנתונה. גרסה SQuAD 2.0 הרחיבה את האתגר בהכנסת שאלות שאין
להן תשובה בפסקה, ובכך מחייבת את המודל לזהות מתי עליו להימנע ממתן תשובה.

מודלים מבוססי טרנספורמר כגון BERT ו-GPT-4 הניחו תשתית חדשה לביצועים על
ערכות אלו. מודל BERT הוכיח כי כוונון עדין (fine-tuning) של מודל מאומן
מראש על משימה ספציפית מניב שיפור ניכר בביצועים בהשוואה לגישות קודמות.
מדדי ההערכה המשמשים על ערכות GLUE ו-SuperGLUE כוללים בדרך כלל שילוב של
דיוק (Accuracy), F1 ומתאם מת'יוס (Matthews Correlation Coefficient,
MCC).

\hypertarget{ux5d4ux5e2ux5e8ux5dbux5ea-ux5deux5d5ux5d3ux5dcux5d9ux5dd-ux5e2ux5d1ux5e8ux5d9ux5d9ux5dd}{%
\subsection{הערכת מודלים
עבריים}\label{ux5d4ux5e2ux5e8ux5dbux5ea-ux5deux5d5ux5d3ux5dcux5d9ux5dd-ux5e2ux5d1ux5e8ux5d9ux5d9ux5dd}}

הערכת מודלי עיבוד שפה טבעית לעברית מציבה אתגרים ייחודיים הנובעים
ממאפייני השפה. העברית היא שפה שמית בעלת מורפולוגיה עשירה, ומילה אחת
עשויה לשאת בתוכה מידע דקדוקי רב הכולל מין, מספר, יחס וכינויי שייכות.
לפיכך, הערכת מודלים עבריים מחייבת מדדים המתחשבים בפירוק מורפולוגי
(morphological disambiguation) ובניתוח תחבירי ייעודי.

מודל AlephBERT, שפותח לצורך עיבוד שפה עברית, הוערך על סדרת משימות עבריות
הכוללות תיוג חלקי דיבר (POS tagging), זיהוי ישויות שמות, וניתוח תחבירי
תלותי. ביצועי המודל הושוו לקווי בסיס קודמים תוך שימוש במדד F1 לרמת
האסימון.

אתגר נוסף בהערכה הוא היעדר ערכות הערכה גדולות ומגוונות בעברית בהשוואה
לאנגלית. מאמצים נמשכים ליצירת ערכות נתונים ייחודיות לעברית, המשמשות הן
לאימון והן להערכה. בנוסף, מולטי-לינגואליות של מודלים כגון mBERT ו-XLM-R
נבחנת באמצעות העברת משימות (zero-shot transfer) מאנגלית לעברית, ומדדי
ההערכה מראים כי מודלים ייעודיים לעברית עולים על ביצועי המודלים
הרב-לשוניים הכלליים.
